\chapter{Introduction\label{cha:chapter1}}

The surge of popularity of open data repositories, scientific data lakes, and enterprise data catalogs (among others) has presented us with the challenge of discovering datasets, in cases where access to the raw data is not practical, or even restricted in some cases. Depending on the data vendor, such as one that may distribute privacy-sensitive information, we are only provided with a statistical data summary in the form of the metadata. Distribution-aware dataset search~\cite{behme_2024} addresses this by developing a way for queries to be done over statistical properties of datasets, instead of the aforementioned raw data. It does so by using percentile predicates in the form \emph{``find all columns where the $p$-th percentile exceeds threshold $t$''}.



Fainder~\cite{behme_2024} introduces an index structure for this exact task. A collection of histograms is created using K-means clustering, it then aligns each histogram that is attached to a cluster onto a shared bin grid, it then precomputes cumulative density arrays that enable us to answer percentile predicates via binary search, instead of needing to traverse histograms fully. On the GitTables benchmark of five million histograms, Fainder achieves more than two orders of magnitude speedup over a na\"{i}ve profile-scan baseline while reducing memory footprint from over 1\,TB to 3.2\,GB.

This thesis is about accelerating Fainder. More explicitly, the goal is to re-implement Fainder's query phase in Rust and to understand what actually binds its performance on multi-core hardware used in modern machines. The investigative portion is the crux of this thesis: it uncovers four major binding ceilings on the Sapphire Rapids Xeon, and shows a pattern of \emph{negative composability} across five different optimisation axes. Lastly, we also develop a pre-flight ceiling-identification methodology from that pattern. All these findings rest on the access pattern that Fainder is implemented against, which is per-cluster binary search over respective density arrays. Fainder here is the empirical measuring stick. The findings throughout this thesis also apply to any other percentile-predicate query engine with that type of access pattern, which makes the methodology here portable beyond its codebase.

\section{Motivation\label{sec:moti}}

Fainder's original implementation~\cite{behme_2024} was implemented in Python, and uses NumPy. Using Python, means that Fainder's code is subject to the Global Interpreter Lock (GIL)~\cite{cpython_gil} which is responsible for serialising bytecode across threads in a single process. Due to the GIL, multi-threading the cluster-chain loop cannot reduce wall-clock query time. The reference implementation also uses a memory layout that is not suited to Fainder's access pattern, and per-iteration interpreter dispatch over the entire cluster chain compounds onto both.

All of this results in the query time of the reference implementation running in the interpreter, not in the binary-search algorithm it revolves around. Python-based Fainder takes hours ($\sim 2.5$ hours on 10,000 queries) on the five-million-histogram primary workload on a modern two-socket server (§\ref{sec:method:hw}); the GIL flattens the thread sweep so no thread count meaningfully reduces this wall. At a throughput floor of hours per workload, the interactive-exploration use case becomes untenable.


\section{Research Challenge\label{sec:objective}}

The challenge in this thesis is to address and attempt to close the gap between Fainder's algorithmic potential, and the performance limits set by its Python implementation, without sacrificing the correctness guarantees that its existing use depends upon. The investigative challenge is the characterisation of \emph{the bounds that are present once the Python overhead is removed}, which hardware ceilings bind, and under which regime, and how the ablations targeting one ceiling work with another.

To perform this research, we need to perform systematic low-level profiling to establish and identify the dominant bottlenecks of performance, design a replacement query engine which is able to make use of cache-conscious memory layouts and the parallelism provided by the hardware. Lastly, we would also need to validate that each of these implementations produces \emph{bitwise-identical} results to the original reference implementation across all workloads. It is also vital that the solution of this research integrates into the current Python implementation and ecosystem, since the rest of Fainder's core (such as clustering and data-loading) is efficiently implemented in Python. The rebinning phase of the index construction has been ported to Rust in Chapter~\ref{cha:chapter3}, but the K-means and histogram construction remain implemented in Python.

\section{Scope}
\textbf{In scope.} The thesis focuses on re-engineering Fainder's query-execution phase at a low system level and empirically evaluating hardware-oriented optimisations. The baseline implementation is first profiled to identify the CPU and memory bottlenecks during percentile evaluation and binary search. Based on the profiling results, we introduce targeted optimisations, including cache-local data layouts for the precomputed index, SIMD vectorisation of cumulative-density comparisons, and multicore parallelisation of independent cluster queries, and assess the performance impact of each, while maintaining exact correctness of query results. Histogram alignment belongs to the index-construction phase and is outside the primary optimisation focus; the alignment (rebinning) kernel was nevertheless ported to Rust, and its construction-time impact is evaluated separately (§\ref{sec:exp:rebinning}).

\textbf{Out of scope.} The thesis does not modify Fainder's algorithmic foundations, mathematical model, or percentile-predicate semantics. It does not explore GPU acceleration, distributed computing, or cloud-based scaling approaches. User-interface design, visualisation, and non-performance-related extensions are likewise excluded. The work is confined to single-node, CPU-level optimisation; all improvements remain transparent to Fainder's logical behaviour, so query results and accuracy remain identical to the reference implementation.

\section{Four Performance Ceilings on Sapphire Rapids\label{sec:ceilings}}

The main empirical finding of this thesis (as unearthed by our ablations across more than twenty optimisation axes) is that distribution-aware search on a Sapphire Rapids server (§\ref{sec:method:hw}) is subject to binding by \textbf{four performance ceilings}, each having a distinct binding under a different regime, addressed with a class of intervention.

\begin{description}
    \item[(i) Single-thread compute / L1 latency.] When we have a thread count of  $ t \leq 8$, the performance wall is the latency chain of dependent loads in the binary search. No latency-side intervention moves the wall beyond noise; SIMD is only useful in a regime where no other ceiling has already collected the wall budget (§\ref{sec:exp:c1}).

    \item[(ii) Shared L3 capacity.] At thread counts in the intermediate range the working sets of each thread begin to overlap in the shared L3. Optimisations that reduce the per-cluster footprint contribute inside the composed bundle, while optimisations that expand the footprint are at a loss in this regime (§\ref{sec:exp:c2}).

    \item[(iii) On-socket memory-traffic pressure.]
    At thread counts in the higher range our ceiling becomes an anti-scaling wall signature, characterised as aggregate memory-hierarchy pressure. Optimisations that reduce the bytes moved per comparison or per emitted result address this ceiling (§\ref{sec:exp:c3}).

    \item[(iv) Cross-socket UPI interconnect.]
    At thread counts which require both sockets, the interconnect between the two NUMA nodes becomes a fourth ceiling. Single-socket placement removes part of the wall time, while naively interleaving memory across both sockets regresses (§\ref{sec:exp:c4}).
    
\end{description}

These four ceilings are what Chapter~\ref{cha:chapter4} works through with evidence: each optimisation axis we investigate was chosen because one of these ceilings predicted it should matter, so the ablation chapter reads as \emph{evidence for each ceiling}.

\section{Negative-Composability\label{sec:negcomp}}

The characterisation of four ceilings is the thesis's principal finding on the \emph{empirical} side. On the other hand, the \emph{structural} finding is that the ceilings are not independent: they interact in a way that makes the composition of optimisations across ceilings non-trivial. Five independent ablation pairs, each on a different hardware axis, exhibit the same shape: a coarser optimisation has already addressed the binding ceiling, and the finer optimisation's overhead ends with a net regression rather than diminishing returns. Each case has an argument for positive composition, which is then falsified by the \texttt{perf} counters.

Pre-flight ceiling-identification is the methodological consequence of this pattern: probe the binding ceiling at the \emph{target composition} before committing to implement a finer optimisation. We demonstrate the five instances in §\ref{sec:exp:negcomp:instances}, and synthesise the pattern in §\ref{sec:exp:negcomp}.

\section{Dispatch-on-Regime\label{sec:dispatch}} The main findings in this Thesis are the four performance ceilings, and negative composability. The dispatch policy aggregates these experimental findings into a deployment artefact: it selects a ``best in regime'' combination of build features and deployment flags. On the 18-cell calibration grid it picks the empirical winner in 17 cases and lands within 5\% in all 18, on the out-of-distribution density point it holds all 6 regime assignments (§\ref{sec:exp:dispatch}, §\ref{sec:exp:ood}).


\section{Contributions\label{sec:contributions}}

The thesis's contributions are organised as follows:

\begin{enumerate}[label=(\arabic*),left=0pt]
    \item \textbf{Rust re-implementation of Fainder's query pipeline.} A Rust implementation of Fainder's query engine, preserving the index structure and interface compatibility with the Python API. We validated the engine to produce bitwise-identical results to the Python reference implementation across all workloads (§\ref{sec:correctness}). At the peak-scaling regime, the end-to-end speedup over the Python baseline spans $530$--$598\times$; most of this ratio is Rust avoiding Python's per-result materialisation cost, and the search phase itself converges on the shared hardware ceilings as per-cluster footprint grows (§\ref{sec:exp:e2e}). The Rust rebinning kernel achieves a separate speedup on the index-construction alignment phase (§\ref{sec:exp:rebinning}). The engineering is presented in Chapter~\ref{cha:chapter3} as our main driver for the ablations.

    \item \textbf{Four-ceiling micro-architectural characterisation.} Identification of four distinct performance ceilings via \texttt{perf} counters on the Sapphire Rapids (Xeon 8468H) server. Each of these ceilings binds under a regime that is addressable through a specific intervention. The evidence for the ceilings (Section~\ref{sec:ceilings}) is presented in Chapter~\ref{cha:chapter4}. 
    
    \item \textbf{Negative-composability pattern.} Five independent instances on five different axes of the same compositional shape, leading us to the result of negative composability of optimisation. This shows as a structural property of the access pattern rather than a property of any single ablation (Section~\ref{sec:negcomp}, evidence and mechanism are presented in §\ref{sec:exp:negcomp:instances}, and synthesis is done in §\ref{sec:exp:negcomp}).
    

    \item \textbf{Pre-flight ceiling-identification methodology.} We identify binding ceilings at the \emph{target composition} using data from \texttt{perf}. These help us to avoid implementing an optimisation that is expected to be positive, but is actually negative at the composition baseline. The methodology is presented in §\ref{sec:exp:negcomp:preflight}.

    \item \textbf{Dispatch-on-regime deployment recipe.} An automatic dispatch that aggregates the experimental findings into a selection of build features and deployment flags. Validated 17/18 exactly and 18/18 within 5\% on the calibration data scales (Section~\ref{sec:dispatch}, derived in §\ref{sec:exp:dispatch}, OOD validation in §\ref{sec:exp:ood}).
\end{enumerate}

The engineering artefacts (1 and 5) are the empirical staging ground for the scientific and methodological findings (2, 3, and 4). We would not have been able to identify the four ceilings without the Rust engine, nor identify the negative composability pattern without the ablations. The pattern itself is what generalises beyond this codebase.

\section{Thesis Structure}

The remainder of this thesis is organised as follows:
\begin{enumerate}
    \item Chapter~\ref{cha:chapter2} provides the scientific background on distribution-aware dataset search, the Fainder index structure, and the hardware-conscious optimisation techniques applied in this work.
    \item Chapter~\ref{cha:chapter3} documents the experimental setup (hardware, software, datasets, queries), the Rust implementation architecture, the measurement protocol, and the a-priori hypotheses each experiment tests.
    \item Chapter~\ref{cha:chapter4} reports the experimental evaluation as evidence for each of the four ceilings identified in Section~\ref{sec:ceilings}, with the negative-composability pattern in its closing section.
    \item Chapter~\ref{cha:chapter5} surveys related work, with particular attention to the composability of optimisations in classical query optimisation and to the joint design of algorithms and hardware.
    \item Chapter~\ref{cha:chapter6} concludes and outlines directions for future work, including the experiments most likely to confirm or extend the structural finding.
    \item Appendix~\ref{app:reproduction} provides the reproduction and implementation details.
\end{enumerate}


